\documentclass[12pt]{article}

\usepackage{sbc-template}
\usepackage[T1]{fontenc}
\usepackage[utf8]{inputenc}
\usepackage[brazil]{babel}
\usepackage{graphicx,url}
\usepackage{amsmath,amssymb,amsthm}
\usepackage{booktabs}
\usepackage{algorithm}
\usepackage{algpseudocode}
\usepackage{tikz}
\usetikzlibrary{fit}

\sloppy

% pseudocódigo em português
\floatname{algorithm}{Algoritmo}
\algrenewcommand\algorithmicrequire{\textbf{Entrada:}}
\algrenewcommand\algorithmicensure{\textbf{Saída:}}
\algrenewcommand\algorithmicwhile{\textbf{enquanto}}
\algrenewcommand\algorithmicdo{\textbf{faça}}
\algrenewcommand\algorithmicfor{\textbf{para}}
\algrenewcommand\algorithmicforall{\textbf{para cada}}
\algrenewcommand\algorithmicif{\textbf{se}}
\algrenewcommand\algorithmicthen{\textbf{então}}
\algrenewcommand\algorithmicend{\textbf{fim}}
\algrenewcommand\algorithmicreturn{\textbf{devolva}}

\newtheorem{teorema}{Teorema}
\newtheorem{corolario}{Corolário}

\newcommand{\chiP}{\chi_P}

\title{Uma Heurística Gulosa por Classes de Cor para o\\
  Problema de Coloração Seletiva de Grafos}

\author{Carlos Alberto de Lima Neto, Guilherme Augusto Pondofe Santana}

% TODO: confirmar o departamento e preencher os e-mails.
\address{Departamento de Informática e Matemática Aplicada (DIMAp)\\
  Universidade Federal do Rio Grande do Norte (UFRN) -- Natal -- RN -- Brasil
}

\begin{document}

\maketitle

\begin{resumo}
  No Problema de Coloração Seletiva de Grafos, os vértices de um grafo estão
  divididos em $k$ grupos, e é preciso escolher exatamente um vértice de cada
  grupo e colorir os escolhidos com o menor número de cores. O problema
  generaliza a coloração de vértices e, por isso, é NP-difícil e difícil de
  aproximar. Implementamos uma heurística construtiva gulosa que monta uma
  classe de cor por vez, provamos que ela roda em tempo
  $O(n\log n + cn + m)$ e usa no máximo $1+\max_i\min_{v\in V_i} d_i(v)$
  cores, e a avaliamos em 280 instâncias da literatura. Todas as soluções foram
  válidas, a maior instância levou 3,4~ms, e o gap médio em relação aos valores
  ótimos ou melhores conhecidos foi de 19\% nas instâncias nsfnet e ring e de
  43\% nas instâncias random.
\end{resumo}

% ---------------------------------------------------------------------------
\section{Introdução}

Problemas de coloração de grafos modelam situações em que recursos precisam
ser atribuídos a elementos que entram em conflito. Na coloração seletiva, além
de colorir, é preciso escolher \emph{quais} elementos serão coloridos: os
vértices estão organizados em grupos e apenas um vértice de cada grupo entra na
solução. O problema foi proposto por \cite{li2000}, com o nome de
\emph{Partition Coloring Problem} (PCP), para modelar o roteamento e a
atribuição de comprimentos de onda em redes ópticas: cada grupo reúne as rotas
possíveis de uma conexão, duas rotas que compartilham um enlace são adjacentes
e cada cor é um comprimento de onda. Outras aplicações, como problemas de
escalonamento e de alocação, são discutidas por \cite{demange2015}.

Uma forma de ver o problema é a seguinte. Há tarefas divididas em categorias, e
de cada categoria basta executar uma tarefa. Duas tarefas são incompatíveis
quando acontecem no mesmo horário. Cada trabalhador executa um conjunto de
tarefas compatíveis entre si. Queremos escolher uma tarefa de cada categoria e
contratar o menor número de trabalhadores capaz de executá-las.

Este trabalho apresenta uma heurística construtiva para o problema. As
contribuições são: (i) a implementação de uma heurística gulosa que monta uma
classe de cor por vez, inspirada na heurística inicial de \cite{furini2018}, e
de um método de duas fases usado como referência de comparação; (ii) a análise
de complexidade de tempo da heurística e um limite superior para o número de
cores que ela usa; e (iii) a aplicação dos dois métodos às 280 instâncias
clássicas da literatura, com comparação aos valores ótimos publicados e a
heurísticas construtivas da literatura.

O restante do texto está organizado assim. A Seção~\ref{sec:problema} define o
problema e a Seção~\ref{sec:complexidade} discute a sua complexidade. A
Seção~\ref{sec:literatura} resume os métodos já publicados e a
Seção~\ref{sec:instancias} apresenta os bancos de instâncias e links úteis. A
Seção~\ref{sec:algoritmo} descreve o algoritmo e a sua implementação, analisados
na Seção~\ref{sec:analise}. A Seção~\ref{sec:experimentos} traz os experimentos
e a Seção~\ref{sec:conclusao} conclui o trabalho.

% ---------------------------------------------------------------------------
\section{Descrição do problema}\label{sec:problema}

Seja $G=(V,E)$ um grafo simples com $n=|V|$ vértices e $m=|E|$ arestas, e seja
$\mathcal{P}=\{V_1,\dots,V_k\}$ uma partição de $V$ em $k$ grupos não vazios.
Uma \emph{seleção} é um conjunto $S\subseteq V$ com $|S\cap V_i|=1$ para todo
$i$, isto é, um representante por grupo. Uma \emph{coloração própria} de $S$ é
uma função $f:S\to\{1,\dots,c\}$ tal que $f(u)\neq f(v)$ para toda aresta
$uv\in E$ com $u,v\in S$. O problema consiste em encontrar uma seleção $S$ e
uma coloração própria de $G[S]$ com o menor número possível de cores. Esse
mínimo é o \emph{número cromático de partição} de $G$, denotado por $\chiP(G)$:
\[
  \chiP(G)=\min\{\chi(G[S]) : S \text{ é uma seleção}\},
\]
onde $\chi$ é o número cromático usual. Arestas entre vértices de um mesmo grupo
nunca aparecem em $G[S]$ e podem ser ignoradas. Para um vértice $v\in V_i$,
escrevemos $d_i(v)=|N(v)\setminus V_i|$ para o número de vizinhos de $v$ fora
do seu grupo.

A Figura~\ref{fig:toy} mostra a instância de exemplo usada nos testes do
projeto. Os vértices ímpares formam um $K_4$, de modo que escolher o primeiro
vértice de cada grupo exige quatro cores. Escolhendo os vértices
$\{2,3,6,8\}$, duas cores bastam, e esse é o ótimo, pois nenhuma seleção é um
conjunto independente.

\begin{figure}[ht]
  \centering
  \begin{tikzpicture}[
      v/.style={circle,draw,minimum size=7mm,inner sep=0pt,font=\small},
      c1/.style={v,fill=black!20},
      c2/.style={v,fill=black,text=white},
      g/.style={draw,dashed,rounded corners,inner sep=2pt}]
    \node[v]  (1) at (0,1.6) {1};  \node[c1] (2) at (0,0) {2};
    \node[c1] (3) at (2.4,1.6) {3}; \node[v]  (4) at (2.4,0) {4};
    \node[v]  (5) at (4.8,1.6) {5}; \node[c1] (6) at (4.8,0) {6};
    \node[v]  (7) at (7.2,1.6) {7}; \node[c2] (8) at (7.2,0) {8};
    \foreach \a/\b in {1/3,3/5,5/7,2/4,6/8,2/5,4/7,1/6,3/8}
      \draw (\a) -- (\b);
    \draw (1) to[bend left=30] (5);
    \draw (3) to[bend left=30] (7);
    \draw (1) to[bend left=40] (7);
    \foreach \a/\b/\i in {1/2/1,3/4/2,5/6/3,7/8/4}
      \node[g,fit=(\a)(\b),label=below:{$V_\i$}] {};
  \end{tikzpicture}
  \caption{Instância de exemplo com $n=8$, $m=12$ e $k=4$. Os vértices
    preenchidos formam a solução ótima encontrada pela heurística: 2, 3 e 6 com
    a cor 1 (cinza) e 8 com a cor 2 (preto).}
  \label{fig:toy}
\end{figure}

% ---------------------------------------------------------------------------
\section{Complexidade do problema}\label{sec:complexidade}

A versão de decisão do problema (``existe seleção colorível com no máximo $c$
cores?'') pertence a NP: uma seleção e uma coloração formam um certificado
verificável em tempo $O(n+m)$. O problema é NP-difícil porque contém a coloração
de vértices como caso particular: quando todo grupo tem um único vértice, a
única seleção é $S=V$ e $\chiP(G)=\chi(G)$. Como decidir se $\chi(G)\le 3$ é
NP-completo \cite{karp1972,garey1979}, decidir se $\chiP(G)\le 3$ também é
\cite{li2000}. A coloração seletiva continua NP-difícil mesmo em classes de
grafos nas quais a coloração usual é fácil, conforme mostrado por
\cite{demange2014}.

A mesma redução vale para a aproximação. \cite{zuckerman2007} mostrou que, para
todo $\varepsilon>0$, não existe algoritmo polinomial que aproxime $\chi(G)$
dentro de um fator $n^{1-\varepsilon}$, a menos que P${}={}$NP. Na redução
acima, $k=n$, de modo que o PCP também não admite aproximação dentro de
$k^{1-\varepsilon}$ sob a mesma hipótese. Em particular, não se espera algoritmo
aproximativo com razão constante para o problema, o que justifica o uso de
heurísticas.

% ---------------------------------------------------------------------------
\section{Heurísticas e algoritmos da literatura}\label{sec:literatura}

\paragraph{Heurísticas construtivas.} \cite{li2000} adaptaram ao PCP três
heurísticas clássicas de coloração: \emph{Largest-First}
\cite{welsh1967}, \emph{Smallest-Last} e \emph{Color-Degree}, esta última
próxima do DSATUR \cite{brelaz1979}. Cada uma tem uma versão de duas fases
(\emph{Twostep}: primeiro escolhe os representantes, depois colore) e uma de
uma fase (\emph{Onestep}: escolhe e colore ao mesmo tempo). A variante
Onestep-CD, que a cada passo colore o vértice de maior grau de cor entre os
grupos ainda não atendidos, foi a melhor delas e serviu de solução inicial para
vários trabalhos posteriores. \cite{fritz2013} propôs uma construção híbrida
que combina Onestep-CD com a heurística DANGER.

\paragraph{Metaheurísticas.} \cite{noronha2006} partem da solução de
Onestep-CD e reduzem o número de cores com busca tabu, superando as heurísticas
construtivas. \cite{pop2013} propuseram um algoritmo memético que combina
operadores genéticos e busca local. \cite{furini2018} propuseram uma busca
adaptativa em vizinhança grande (ALNS) com busca local, cuja solução inicial é
gerada por uma heurística gulosa que constrói um conjunto independente por vez,
a mesma ideia do método implementado neste trabalho.

\paragraph{Algoritmos exatos.} \cite{frota2010} apresentaram um algoritmo
branch-and-cut baseado na formulação por representantes.
\cite{hoshino2011} propuseram um branch-and-price sobre uma reformulação de
Dantzig-Wolfe da mesma formulação. \cite{furini2018} propuseram um novo
branch-and-price, com subproblema de conjunto independente de peso máximo, que
resolveu pela primeira vez várias instâncias em aberto. Os valores ótimos
publicados nesse trabalho são a principal referência dos nossos experimentos.

\paragraph{Algoritmos aproximativos.} Pelo resultado de inaproximabilidade da
Seção~\ref{sec:complexidade}, não se conhecem algoritmos aproximativos com
razão constante para o PCP, e não se espera que existam. Os métodos práticos
são heurísticos ou exatos, como os citados acima.

% ---------------------------------------------------------------------------
\section{Bancos de instâncias e links}\label{sec:instancias}

Usamos o banco clássico de instâncias do PCP, distribuído pelo grupo de C. C.
Ribeiro e mantido no repositório de \cite{santini2017}, que acompanha
\cite{furini2018}. A Tabela~\ref{tab:instancias} resume os três conjuntos.

\begin{itemize}
  \item \textbf{random} \cite{frota2010}: grafos aleatórios com grupos de dois
    vértices. Há um conjunto com densidade 0,5 e $n$ de 20 a 120, e outro com
    $n=90$ e densidade de 0,1 a 0,9, com cinco instâncias por combinação.
  \item \textbf{nsfnet} \cite{li2000}: roteamento na rede NSF (14 nós e 21
    enlaces). Cada conexão pedida com probabilidade $p\in\{0{,}1,\dots,1{,}0\}$
    gera um grupo com até duas rotas, e duas rotas são adjacentes quando
    compartilham um enlace.
  \item \textbf{ring} \cite{frota2010}: o mesmo processo em redes em anel com
    10, 15 e 20 nós. Os cinco arquivos \texttt{ring\_n\{10,15,20\}p1.0s\{1..5\}}
    de cada tamanho são cópias idênticas, como notam \cite{furini2018}; eles
    foram mantidos para preservar os nomes originais.
\end{itemize}

\begin{table}[ht]
  \centering
  \caption{Conjuntos de instâncias usados nos experimentos.}
  \label{tab:instancias}
  \begin{tabular}{lrcccl}
    \toprule
    Conjunto & Arquivos & $n$ & $k$ & Densidade & Origem \\
    \midrule
    random & 80  & 20--120 & 10--60  & 0,10--0,91 & \cite{frota2010} \\
    nsfnet & 50  & 16--257 & 12--182 & 0,09--0,13 & \cite{li2000} \\
    ring   & 150 & 8--760  & 4--380  & 0,36--0,41 & \cite{frota2010} \\
    \bottomrule
  \end{tabular}
\end{table}

As instâncias foram convertidas para um formato textual simples, com vértices e
grupos numerados a partir de 1: uma linha \texttt{p edge $n$ $m$ $k$}, uma linha
\texttt{c $i$ $v_1$ $v_2$ \dots} por grupo e uma linha \texttt{e $u$ $v$} por
aresta. Como existem instâncias da literatura com tamanho significativo (até
760 vértices e 117 mil arestas), não foi necessário criar um gerador.

Links importantes:
\begin{itemize}
  \item Repositório deste trabalho (código, instâncias e resultados):
    \url{https://github.com/pondof/trab1PAA}.
  \item Instâncias e código de \cite{furini2018}:
    \url{https://github.com/alberto-santini/selective-graph-colouring}.
  \item Página do grupo de C. C. Ribeiro sobre o PCP:
    \url{http://www2.ic.uff.br/~celso/grupo/pcp.htm}.
\end{itemize}

% ---------------------------------------------------------------------------
\section{Algoritmo desenvolvido}\label{sec:algoritmo}

\subsection{Guloso por classes de cor}

A heurística principal, chamada \texttt{guloso\_classes}, contrata um
trabalhador (cor) por vez e dá a ele o maior número possível de tarefas
compatíveis, no máximo uma por grupo. Os vértices são percorridos em ordem
crescente de grau, com desempate pelo menor identificador: vértices de grau
baixo conflitam com poucos outros e tendem a deixar espaço para mais vértices
na mesma cor. Um vértice entra na cor atual se o seu grupo ainda não foi
atendido e se ele não é vizinho de nenhum vértice já colocado nessa cor. Quando
a varredura termina, a cor está fechada e uma nova cor é aberta. O processo
para quando todos os grupos têm representante. O Algoritmo~\ref{alg:guloso}
resume o método, baseado na heurística gulosa do código de \cite{santini2017}.

\begin{algorithm}[ht]
  \caption{Guloso por classes de cor}\label{alg:guloso}
  \begin{algorithmic}[1]
    \Require grafo $G=(V,E)$ e partição $\{V_1,\dots,V_k\}$
    \Ensure representante $r_i$ e cor $f_i$ de cada grupo $V_i$
    \State $L \gets$ vértices de $V$ em ordem crescente de $(\deg(v), v)$
    \State $P \gets \{1,\dots,k\}$ \Comment{grupos pendentes}
    \State $c \gets 0$
    \While{$P \neq \emptyset$}
      \State $c \gets c+1$; \quad $B \gets \emptyset$ \Comment{$B$: vizinhos da cor $c$}
      \ForAll{$v \in L$}
        \State $i \gets$ grupo de $v$
        \If{$i \in P$ \textbf{e} $v \notin B$}
          \State $r_i \gets v$; \quad $f_i \gets c$; \quad $P \gets P\setminus\{i\}$
          \State $B \gets B \cup N(v)$
        \EndIf
      \EndFor
    \EndWhile
    \State \Return $(r, f)$
  \end{algorithmic}
\end{algorithm}

Na instância da Figura~\ref{fig:toy}, a ordem é $2,4,6,8,1,3,5,7$ (os pares
têm grau 2 e os ímpares, grau 4). Na cor 1 entram 2, 6 e 3; os vértices 4, 8 e
7 ficam bloqueados por vizinhos já escolhidos. Na cor 2 entra o vértice 8,
atendendo o último grupo. A solução usa duas cores, que é o ótimo.

\subsection{Método de duas fases (referência)}

Para ter uma referência de comparação, implementamos também o método
\texttt{duas\_fases}, uma versão Twostep simples: (1) em cada grupo, escolhe o
vértice de menor grau, com desempate pelo menor identificador; (2) colore os
escolhidos com \emph{first-fit} na ordem dos grupos, dando a cada um a menor cor
não usada pelos seus vizinhos já coloridos.

\subsection{Implementação}

Os métodos foram implementados em Python~3, usando apenas a biblioteca padrão.
O grafo é guardado como uma lista de conjuntos de adjacência, o que torna
$O(1)$ (esperado) o teste de adjacência. O código está organizado assim:
\begin{itemize}
  \item \texttt{sgcp.py}: leitura das instâncias, os dois métodos e o
    validador;
  \item \texttt{main.py}: interface de terminal. Um menu permite escolher o
    conjunto e depois a instância; os dois métodos são executados e um
    relatório com os dados da instância, o número de cores, o tempo e a
    validação de cada método e a melhor solução (representante e cor de cada
    grupo e tarefas de cada trabalhador) é mostrado na tela e salvo em
    \texttt{resultados/};
  \item \texttt{experimentos.py}: executa os dois métodos em todas as
    instâncias e gera as tabelas de resultados;
  \item \texttt{testes.py}: testes automatizados com a instância de exemplo.
\end{itemize}
Toda solução passa por um validador antes de ser mostrada ou salva: ele confere
que cada grupo tem um representante que pertence ao grupo e que nenhuma aresta
liga dois representantes da mesma cor. Para executar, basta, na pasta do
projeto, rodar \texttt{python main.py} (menu), \texttt{python main.py
instancias/exemplo/toy.txt} (instância direta) ou \texttt{python
experimentos.py} (todas as instâncias). Em alguns sistemas o comando é
\texttt{python3}.

% ---------------------------------------------------------------------------
\section{Análise do algoritmo}\label{sec:analise}

\subsection{Complexidade de tempo}

Seja $c$ o número de cores usado. A ordenação inicial custa $O(n\log n)$. Cada
iteração do laço externo percorre os $n$ vértices de $L$, com testes de
pertinência $O(1)$, e há $c$ iterações, totalizando $O(cn)$. A atualização
$B\gets B\cup N(v)$ é executada uma única vez por grupo, quando o seu
representante é escolhido, e custa $O(\deg(v))$; somando, o custo é no máximo
$\sum_{v}\deg(v)=2m$. O tempo total é, portanto,
\[
  O(n\log n + cn + m).
\]
Como $c\le k\le n$ (Corolário~\ref{cor:k}), isso é $O(n^2+m)$ no pior caso, e
$O(n\log n + m)$ quando o número de cores é limitado por uma constante. A
memória usada é $O(n+m)$. O método de duas fases custa $O(n)$ na seleção e
$O(k^2)$ no first-fit, que compara cada grupo com os anteriores, totalizando
$O(n+k^2)$. O validador custa $O(k^2)$, e a leitura da instância, $O(n+m)$.

\subsection{Limite para o número de cores}

\begin{teorema}\label{teo:limite}
  O Algoritmo~\ref{alg:guloso} usa no máximo
  \[
    c \;\le\; 1+\max_{1\le i\le k}\ \min_{v\in V_i} d_i(v)
  \]
  cores, onde $d_i(v)$ é o número de vizinhos de $v$ fora do grupo $V_i$.
\end{teorema}

\begin{proof}
  Seja $V_j$ o grupo que recebeu a última cor, $c$. Considere uma cor
  $\ell<c$. Durante a construção da cor $\ell$, o grupo $j$ esteve pendente o
  tempo todo, e todos os vértices de $V_j$ foram examinados sem serem
  escolhidos. Pela condição da linha 8, cada $v\in V_j$ estava em $B$ quando
  foi examinado, isto é, era vizinho de algum vértice $u_\ell$ já colocado na
  cor $\ell$. Esse $u_\ell$ não pertence a $V_j$, pois nenhum vértice de $V_j$
  foi escolhido antes da cor $c$. Como cada vértice recebe no máximo uma cor,
  os vértices $u_1,\dots,u_{c-1}$ são distintos. Logo, todo $v\in V_j$ tem pelo
  menos $c-1$ vizinhos fora de $V_j$, ou seja,
  $c-1\le \min_{v\in V_j} d_j(v)\le \max_i \min_{v\in V_i} d_i(v)$.
\end{proof}

\begin{corolario}\label{cor:k}
  O algoritmo termina e usa no máximo $\min\{k,\ \Delta+1\}$ cores, onde
  $\Delta$ é o grau máximo de $G$.
\end{corolario}

\begin{proof}
  No início de cada cor, $B=\emptyset$; o primeiro vértice de $L$ cujo grupo
  está pendente é sempre escolhido. Assim, cada cor atende pelo menos um grupo,
  o laço termina e $c\le k$. O limite $\Delta+1$ segue do
  Teorema~\ref{teo:limite}, pois $d_i(v)\le\deg(v)\le\Delta$.
\end{proof}

Como $\chiP(G)\ge 1$, o Teorema~\ref{teo:limite} dá uma razão de aproximação
de no máximo $1+\max_i\min_{v\in V_i} d_i(v)$, que depende da instância e não é
constante. Pela Seção~\ref{sec:complexidade}, nenhuma heurística polinomial
pode garantir razão $k^{1-\varepsilon}$, a menos que P${}={}$NP; por isso a
qualidade do método é avaliada experimentalmente na próxima seção.

% ---------------------------------------------------------------------------
\section{Experimentos}\label{sec:experimentos}

\subsection{Metodologia}

Os experimentos foram executados em um notebook com processador Intel Core
i7-1365U, 32~GB de memória e Windows~11, usando CPython~3.13 em uma única
\emph{thread}. Os dois métodos são determinísticos e foram executados uma vez
por instância. O tempo, medido com \texttt{time.perf\_counter}, inclui apenas o
algoritmo, sem a leitura do arquivo. Todas as soluções dos dois métodos, nas
280 instâncias, passaram pelo validador.

Os valores de referência foram retirados das Tabelas~2 e~3 de
\cite{furini2018}, que trazem os resultados de 175 instâncias distintas: 56
random, 32 nsfnet e 87 ring. Quando o limite superior de \cite{hoshino2011},
reproduzido na Tabela~3, era melhor, ele foi usado. Com as cópias idênticas de
\texttt{ring\_n*p1.0}, esses valores cobrem 187 dos nossos arquivos. Em 180
deles o valor é ótimo comprovado; nos outros 7 (\texttt{n90p2t2s4},
\texttt{n90p3t2s4} e as cinco cópias de \texttt{ring\_n20p1.0s1}) é o melhor
valor conhecido. Para conferir que as instâncias são as mesmas, comparamos o
tamanho dos grafos: $n$ coincide em todos os casos, e o $|E|$ informado no
artigo, que conta as arestas depois de transformar cada grupo em uma clique,
coincide em 174 das 175 instâncias. A exceção, \texttt{n90p2t2s2}, tem o mesmo
$|E|$ que \texttt{n90p2t2s4} no artigo, o que sugere um erro de digitação. O
gap de um método é $100\cdot(c-c^*)/c^*$, onde $c^*$ é o valor de referência.

\subsection{Resultados}

A Tabela~\ref{tab:resumo} resume os resultados por conjunto. A coluna
``Melhor/igual/pior'' conta as instâncias em que o guloso usou menos, o mesmo
número ou mais cores que o método de duas fases. Os gaps são médias sobre as
instâncias com valor de referência. A tabela completa, instância por instância,
está no arquivo \texttt{resultados/tabela.md} do repositório.

\begin{table}[ht]
  \centering
  \caption{Resumo por conjunto: médias de cores e tempo sobre todas as
    instâncias; gaps sobre as instâncias com referência.}
  \label{tab:resumo}
  \footnotesize\setlength{\tabcolsep}{4pt}
  \begin{tabular}{lrrrcrrrr}
    \toprule
    & & \multicolumn{2}{c}{Cores (média)} & Melhor/ & Tempo & Com & \multicolumn{2}{c}{Gap médio (\%)} \\
    \cmidrule(lr){3-4}\cmidrule(lr){8-9}
    Conjunto & Inst. & Duas f. & Guloso & igual/pior & (ms) & ref. & Duas f. & Guloso \\
    \midrule
    nsfnet & 50  & 10,04 & 9,66  & 21/24/5  & 0,18 & 32 & 23,81 & 19,17 \\
    random & 80  & 10,56 & 9,74  & 53/20/7  & 0,13 & 56 & 57,23 & 43,30 \\
    ring   & 150 & 22,63 & 22,70 & 49/40/61 & 0,67 & 99 & 18,84 & 18,83 \\
    \bottomrule
  \end{tabular}
\end{table}

A Tabela~\ref{tab:grupos} detalha os resultados por grupo de instâncias,
considerando apenas as instâncias com referência. A coluna $\Delta_{\max}$ é a
maior diferença, em cores, entre o guloso e a referência dentro do grupo.

\begin{table}[ht]
  \centering
  \caption{Resultados por grupo, apenas nas instâncias com referência: médias de
    $n$, do valor de referência ($c^*$) e das cores de cada método, e gap
    médio.}
  \label{tab:grupos}
  \footnotesize\setlength{\tabcolsep}{4pt}
  \begin{tabular}{llrrrrrrrr}
    \toprule
    & & & & & \multicolumn{2}{c}{Cores (média)} & \multicolumn{2}{c}{Gap (\%)} & \\
    \cmidrule(lr){6-7}\cmidrule(lr){8-9}
    Conjunto & Grupo & Inst. & $n$ & $c^*$ & Duas f. & Guloso & Duas f. & Guloso & $\Delta_{\max}$ \\
    \midrule
    random & $n=60$,  $p=0{,}5$ & 1 & 60  & 5,00  & 9,00  & 9,00  & 80,0 & 80,0 & 4 \\
           & $n=70$,  $p=0{,}5$ & 3 & 70  & 6,00  & 9,33  & 8,33  & 55,6 & 38,9 & 3 \\
           & $n=80$,  $p=0{,}5$ & 5 & 80  & 6,00  & 10,40 & 9,20  & 73,3 & 53,3 & 4 \\
           & $n=90$,  $p=0{,}5$ & 5 & 90  & 7,00  & 11,00 & 9,60  & 57,1 & 37,1 & 3 \\
           & $n=100$, $p=0{,}5$ & 5 & 100 & 7,00  & 11,80 & 10,40 & 68,6 & 48,6 & 4 \\
           & $n=120$, $p=0{,}5$ & 5 & 120 & 8,00  & 13,40 & 12,00 & 67,5 & 50,0 & 5 \\
           & $n=90$,  $p=0{,}1$ & 5 & 90  & 2,80  & 4,20  & 3,60  & 53,3 & 30,0 & 1 \\
           & $n=90$,  $p=0{,}2$ & 5 & 90  & 3,80  & 5,60  & 5,20  & 50,0 & 38,3 & 2 \\
           & $n=90$,  $p=0{,}3$ & 5 & 90  & 5,00  & 7,40  & 7,00  & 48,0 & 40,0 & 2 \\
           & $n=90$,  $p=0{,}4$ & 5 & 90  & 5,80  & 8,80  & 8,80  & 52,0 & 52,0 & 4 \\
           & $n=90$,  $p=0{,}6$ & 4 & 90  & 8,25  & 13,25 & 11,75 & 61,1 & 42,7 & 4 \\
           & $n=90$,  $p=0{,}8$ & 4 & 90  & 12,00 & 18,25 & 17,00 & 52,1 & 41,7 & 6 \\
           & $n=90$,  $p=0{,}9$ & 4 & 90  & 16,00 & 22,25 & 21,75 & 39,1 & 35,9 & 7 \\
    \midrule
    nsfnet & $p=0{,}4$ & 3 & 104 & 6,33  & 8,00  & 7,67  & 27,0 & 21,4 & 2 \\
           & $p=0{,}5$ & 4 & 126 & 7,50  & 9,25  & 9,25  & 23,7 & 23,7 & 2 \\
           & $p=0{,}6$ & 5 & 151 & 9,00  & 10,80 & 10,60 & 20,0 & 17,8 & 3 \\
           & $p=0{,}7$ & 5 & 181 & 10,20 & 12,60 & 12,20 & 23,7 & 19,7 & 3 \\
           & $p=0{,}8$ & 5 & 205 & 11,00 & 14,00 & 13,00 & 27,3 & 18,2 & 3 \\
           & $p=0{,}9$ & 5 & 225 & 12,00 & 14,60 & 14,20 & 21,7 & 18,3 & 3 \\
           & $p=1{,}0$ & 5 & 249 & 13,00 & 16,20 & 15,20 & 24,6 & 16,9 & 3 \\
    \midrule
    ring   & 10 nós & 16 & 161 & 12,06 & 14,62 & 14,12 & 21,0 & 17,0 & 3 \\
           & 15 nós & 38 & 287 & 21,50 & 25,03 & 25,53 & 17,4 & 18,4 & 7 \\
           & 20 nós & 45 & 460 & 33,78 & 39,93 & 40,27 & 19,3 & 19,8 & 11 \\
    \bottomrule
  \end{tabular}
\end{table}

Para comparar com heurísticas construtivas publicadas, a
Tabela~\ref{tab:onestep} reúne as médias por conjunto random reportadas por
\cite{fritz2013} para Onestep-CD \cite{li2000} e para a construção híbrida
daquele trabalho, ao lado das médias dos nossos dois métodos nas mesmas
instâncias. Os valores de \cite{fritz2013} vêm de outra implementação, em outra
máquina, por isso comparamos apenas o número de cores.

\begin{table}[ht]
  \centering
  \caption{Média de cores por conjunto random (cinco instâncias cada):
    métodos deste trabalho e heurísticas construtivas reportadas por
    \cite{fritz2013}.}
  \label{tab:onestep}
  \small
  \begin{tabular}{rrrrrr}
    \toprule
    $n$ & $p$ & Duas fases & Guloso & Onestep-CD & Híbrida \\
    \midrule
    20  & 0,5 & 3,8  & 3,4  & 3,6  & 3,4  \\
    40  & 0,5 & 5,8  & 6,0  & 5,6  & 5,0  \\
    60  & 0,5 & 8,2  & 7,6  & 7,4  & 6,6  \\
    70  & 0,5 & 9,2  & 8,4  & 7,6  & 7,0  \\
    80  & 0,5 & 10,4 & 9,2  & 8,6  & 8,0  \\
    90  & 0,5 & 11,0 & 9,6  & 9,4  & 9,4  \\
    100 & 0,5 & 11,8 & 10,4 & 10,2 & 10,0 \\
    120 & 0,5 & 13,4 & 12,0 & 12,0 & 11,0 \\
    90  & 0,1 & 4,2  & 3,6  & 3,0  & 3,0  \\
    90  & 0,2 & 5,6  & 5,2  & 5,0  & 4,2  \\
    90  & 0,3 & 7,4  & 7,0  & 6,2  & 5,8  \\
    90  & 0,4 & 8,8  & 8,8  & 7,8  & 7,6  \\
    90  & 0,6 & 13,2 & 11,8 & 11,2 & 10,4 \\
    90  & 0,7 & 15,2 & 14,0 & 13,6 & 12,6 \\
    90  & 0,8 & 18,6 & 17,4 & 16,4 & 15,6 \\
    90  & 0,9 & 22,4 & 21,4 & 20,0 & 19,8 \\
    \bottomrule
  \end{tabular}
\end{table}

\subsection{Discussão}

\paragraph{Guloso contra duas fases.} O guloso por classes usou, em média,
0,8 cor a menos que o método de duas fases nas instâncias random e 0,4 a menos
nas nsfnet, sendo melhor ou igual em 91\% e 90\% dessas instâncias,
respectivamente. Nas instâncias ring os dois métodos empatam na média
(22,70 contra 22,63 cores, gap de 18,8\% em ambos), e o guloso é pior em 61
das 150 instâncias. Medimos que, em ring e nsfnet, os dois métodos escolhem o
mesmo representante em 97\% e 96\% dos grupos, contra 86\% em random. Ou seja,
nesses conjuntos a diferença está quase toda na coloração: a ordem dos grupos
nos arquivos ring, que segue a enumeração das conexões da rede, parece
favorecer o first-fit tanto quanto a ordem por grau favorece o guloso.

\paragraph{Distância até o ótimo.} O guloso atingiu o valor de referência em
apenas uma das 187 instâncias com referência (\texttt{n90p1t2s1}, com 3
cores), e o método de duas fases, em três. Em média, ele usa 1,9 cor a mais que a referência em nsfnet, 3,0 em
random e 4,8 em ring. O gap relativo é maior em random (43\%) porque essas
instâncias têm poucos grupos e números cromáticos pequenos (de 2 a 16): uma
única cor extra já representa de 6\% a 50\% do ótimo. Nas ring, o gap
relativo é menor (19\%), mas a diferença absoluta cresce com o tamanho da rede
e chega a 11 cores nas maiores instâncias. Para comparação, a ALNS com busca
local de \cite{furini2018}, uma metaheurística que roda por muito mais tempo,
encontra o ótimo em cerca de 158 de 171 instâncias e nunca fica mais de duas
cores acima dele.

\paragraph{Comparação com heurísticas construtivas.} Na
Tabela~\ref{tab:onestep}, Onestep-CD usa de 0,2 a 1,4 cor a menos que o guloso
em 14 dos 16 conjuntos, empata em $n=120$ e perde em $n=20$. A construção
híbrida de \cite{fritz2013} é melhor ainda. A diferença é esperada: Onestep-CD
recalcula, a cada passo, quantas cores diferentes cercam cada vértice (o grau
de cor), enquanto o nosso guloso usa uma ordem fixa, calculada uma única vez a
partir do grau. Essa ordem fixa é o que mantém o método simples e muito rápido.

\paragraph{Tempo.} O guloso levou no máximo 3,4~ms por instância (nas maiores,
as cópias de \texttt{ring\_n20p1.0s1}, com 760 vértices e 117 mil arestas) e
cerca de
0,12~s para as 280 instâncias. O crescimento acompanha o termo $cn+m$ da
análise: nas ring, $m$ domina e o tempo cresce com o número de arestas. Os
métodos exatos de \cite{furini2018} levam de segundos a uma hora nas mesmas
instâncias, o que torna o guloso adequado para gerar soluções iniciais.

% ---------------------------------------------------------------------------
\section{Conclusão}\label{sec:conclusao}

Este trabalho estudou o Problema de Coloração Seletiva de Grafos, NP-difícil e
sem aproximação com razão constante, e implementou uma heurística gulosa que
monta uma classe de cor por vez. Mostramos que ela roda em
$O(n\log n + cn + m)$ e usa no máximo $1+\max_i\min_{v\in V_i} d_i(v)$ cores.
Nas 280 instâncias da literatura, todas as soluções foram válidas, o tempo
ficou abaixo de 4~ms, e o método superou um método de duas fases nas instâncias
random e nsfnet, empatando nas ring. A distância até os valores ótimos
(19\% em nsfnet e ring, 43\% em random) e a comparação com Onestep-CD mostram
que a ordem fixa por grau limita a qualidade das soluções.

Como trabalhos futuros, sugerimos: (i) trocar a ordem fixa por uma ordem
dinâmica, como o grau de cor de Onestep-CD; (ii) aplicar uma busca local que
esvazia a menor classe de cor e tenta reinserir os seus grupos nas demais, como
em \cite{furini2018}; e (iii) executar várias vezes com desempates aleatórios,
ficando com a melhor solução.

\bibliographystyle{sbc}
\bibliography{referencias}

\end{document}
